\section{The Bug, the Fix, and the Experimental Setup}
\label{sec:setup}

\subsection{Template-diff tokenization under \qwen{}}
\label{sec:bug}

The re-implementation stores an agent trajectory as a growing list of chat messages and
tokenizes each new user-role turn, including every tool observation, as the suffix of the
template render of the list with the turn appended, after removing the render of the list
without it. This is exact when the template is prefix-stable, that is, when appending a
message only appends tokens. The \qwen{} template is not prefix-stable. It keeps
\think{}\ldots\texttt{</think>} content only for assistant messages that follow the last plain
user message, so appending an observation deletes every earlier think block from the render.
The render with the observation is then \emph{shorter} than the render without it by the
length of the preceding think block, and the suffix slice starts that many tokens too late.

The consequences are deterministic. The observation always loses its
\texttt{<|im\_start|>user} header and its first $L$ tokens, where $L$ is the length of the
preceding think block. When $L$ exceeds the observation length, the observation vanishes and
the policy sees only the tail of its own prior text. Branch prompts are user-role messages and
are sliced the same way. Applying the parser of \cref{sec:behavior} to the greedy validation
dumps of the earlier study confirms this: none of 913 base-model observations carried a
header, 13--16\% of tool calls in the trained arms received no observation, and 79--82\% of
think blocks were empty across all three arms (\cref{fig:contrast}, grey bars). The policies
still received the tail of most observations, which is why the earlier accuracies were
plausible rather than zero.

\seedoss{}, the paper's model, renders a separate \texttt{reasoning\_content} field and never
parses or strips think tags from the message content, and its template has no
last-user-message logic. The identical code produces an intact user turn for it. The bug is
therefore a property of the Qwen3 template family and of our base-model substitution, not of
the published method.

\paragraph{Fix.} We render every post-prompt turn standalone: the turn is templated behind a
fixed system message and a dummy user message whose token prefix is known and removed, so the
slice no longer depends on history. Tool observations are wrapped in
\toolresp{}\ldots\texttt{</tool\_response>}, the form the \qwen{} template expects for tool
output; with that wrapper the template keeps think blocks in both renders, so server-side
re-templating of a trajectory now agrees with the training tokenization. Eleven tests against
the real \qwen{} tokenizer check header presence, exact token equality with a direct encode,
and think preservation over multi-turn histories.

\subsection{Task, arms, and infrastructure}

We use the authors' released \bcplus{} split (680 training and 150 test instances, 50 each
easy, medium, and hard) with search over the Tevatron \bcplus{} corpus using their precomputed
Qwen3-Embedding-8B index. The tools are \texttt{search}, \texttt{open\_page},
\texttt{finish}, and, for the folding agent, \texttt{branch} and \texttt{return}. Both RL arms
start from \qwen{} with thinking mode on. \foldgrpo{} runs the authors' training script
verbatim; \grpo{} differs by exactly two flags (outcome-only advantages, no process rewards).
\Cref{tab:config} lists the configuration against the paper's.

All judge roles (training outcome reward, \foldgrpo{}'s out-of-scope penalty, and final
grading) are served by \judgemodel{} behind an API shim, so absolute scores are not comparable
to the paper's; every cross-arm comparison uses this one fixed judge. Training ran as one
8$\times$H100 batch job per arm on a shared queue, with a separate 8$\times$A100 node serving
the search index and the judge. The shim re-resolves that node's address on every request, so a
replaced node is picked up without restarting training. Checkpoints were saved every 10 steps
and final evaluations ran as self-contained jobs that host search, judge, and rollout on one
node.

\paragraph{Evaluation protocol.} All scores are measured through the training stack's
validation path at two operating points: greedy decoding (pass@1, $n{=}150$, binomial 95\% CI
about $\pm 7.5$ points at $p \approx 0.3$) and temperature 1.0 with four samples per item
(mean@4, 600 graded trajectories, nominal CI about $\pm 3.7$ points; the four samples per item
are clustered, so the paired CI for an arm difference is about $\pm 3.5$ points). The
temperature-1.0 tier was pre-registered as decisive. Every 10 training steps the trainer also
runs the greedy validation in-process; validations that fell inside an infrastructure gap
(\cref{sec:budget}) were re-run from the uploaded checkpoint.

\begin{table}[t]
\centering
\small
\begin{tabular}{p{4.2cm}p{3.6cm}p{6.6cm}}
\toprule
Setting & Paper & This campaign \\
\midrule
Base model & \seedoss{} & \qwen{}, thinking mode on \\
Judge (all roles) & GPT-5-nano / gpt-4o-mini / gpt-4.1 & \judgemodel{} (local vLLM) \\
Rollout batch / group / PPO mini-batch & 32 / 8 / 128 & 32 / 8 / 128 \\
Learning rate; response tokens; max branches & $10^{-6}$; 32{,}768; 10 & $10^{-6}$; 32{,}768; 10 \\
Observation rendering & template renders history verbatim & standalone turn render, \toolresp{} wrapping (\cref{sec:bug}) \\
Training steps & 50 & 100 planned; 46 (\foldgrpo{}) and 59 (\grpo{}) reward-bearing (\cref{sec:budget}) \\
Compute & --- & $8\times$H100 per arm; $8\times$A100 search + judge node \\
\bottomrule
\end{tabular}
\caption{Configuration versus the paper. Training hyperparameters follow the authors' released
script verbatim; the base model, the judge, the observation rendering, and the realized step
budget differ.}
\label{tab:config}
\end{table}
